\documentclass[12pt,a4paper]{article}

% PAQUETES
\usepackage[utf8]{inputenc}
\usepackage[T1]{fontenc}
\usepackage[spanish]{babel}
\usepackage{amsmath}
\usepackage{geometry}
\usepackage{graphicx}
\usepackage{float}
\usepackage{listings}
\usepackage{xcolor}
\usepackage{setspace}

% CONFIGURACIÓN
\geometry{margin=2.5cm}
\onehalfspacing

\begin{document}

% PORTADA
\begin{titlepage}
    \centering
    
    \vspace*{1.5cm}
    
    {\Large \textbf{UNIVERSIDAD DEL CAUCA}}\\[0.5cm]
    {\large Facultad de Ingeniería Electrónica y Telecomunicaciones}\\[2cm]
    
    {\Huge \textbf{DIGITALIZACIÓN DE SEÑALES}}\\[0.3cm]
    {\Large Conmutación telefónica mediante PCM}\\[2cm]
    
    \textbf{Práctica de laboratorio}\\[1.5cm]
    
    \begin{center}
        \textbf{Estudiante: Dhylan Styven Collazos Caicedo}\\[0.5cm]
        \textbf{Asignatura: Sistemas de Conmutación}\\[0.5cm]
        \textbf{Docente: Carlos Hernan Tobar Arteaga}\\[0.5cm]
        \textbf{Fecha: 11/08/2026}
    \end{center}
    
    \vfill
    
    {\large Popayán, Cauca}\\
    {\large 2026}
\end{titlepage}

% INTRODUCCIÓN
\section{Introducción}

La digitalización de señales permite transformar señales
analógicas, como la voz, en información digital que puede
transmitirse y procesarse eficientemente.

En esta práctica se estudia la modulación por impulsos
codificados (PCM), considerando las etapas de muestreo,
cuantificación y codificación. También se analiza el uso
de companding A-law, la transmisión mediante un canal
ruidoso y la reconstrucción de la señal en el receptor.

A continuación, se presentan las respuestas a las preguntas
planteadas durante el desarrollo de la práctica.

% DESARROLLO
\section{Desarrollo de las preguntas}

\subsection{Frecuencia de muestreo}

\textbf{¿Qué sucede al disminuir la frecuencia de muestreo
de 8000 Hz a 4000 Hz?}

Al disminuir la frecuencia de muestreo se reduce la cantidad
de muestras tomadas por segundo, provocando pérdida de
información de la señal original.

Según el teorema de Nyquist-Shannon, se requiere:

\[
F_s \geq 2f_{\max}
\]

Por lo tanto, con 4000 Hz solamente podrían representarse
frecuencias de hasta 2000 Hz sin aliasing. Esto ocasiona
pérdida de componentes importantes de la voz y disminuye
la claridad del audio reconstruido.

\subsection{Bits por muestra}

\textbf{¿Cómo varía el ruido de cuantización al utilizar
4, 8 o 12 bits?}

La cantidad de bits determina los niveles disponibles para
representar la amplitud de la señal:

\[
L=2^n
\]

Con 4 bits se obtienen 16 niveles, con 8 bits se obtienen
256 y con 12 bits se obtienen 4096 niveles.

Al aumentar la resolución disminuye el error de
cuantización, permitiendo representar la señal con mayor
precisión y mejorar la calidad del audio.

\subsection{Companding A-law vs $\mu$-law}

\textbf{¿Qué diferencias existen en la forma de onda y
la calidad percibida?}

Ambas técnicas utilizan compresión logarítmica para mejorar
la representación de señales de voz.

A-law se utiliza principalmente en Europa, mientras que
$\mu$-law es común en Norteamérica. Aunque ambas reducen
el error de cuantificación en señales de voz, presentan
diferencias en la distribución de los niveles y en la
forma de onda reconstruida.

\subsection{Número de niveles de cuantificación}

\textbf{¿Qué ocurre al utilizar 16, 256 o 1024 niveles?}

Con 16 niveles la señal presenta mayor distorsión debido
a la baja resolución. Con 256 niveles se obtiene una
representación más precisa, mientras que con 1024 niveles
disminuye el error de cuantificación.

Por lo tanto, aumentar los niveles mejora la fidelidad
de la señal, aunque requiere más bits para su transmisión.

\subsection{Relación señal/ruido (SNR)}

\textbf{¿Cuál es el impacto de utilizar 30 dB, 10 dB y 0 dB?}

La relación señal/ruido se expresa mediante:

\[
SNR_{dB}=10\log_{10}\left(\frac{P_s}{P_n}\right)
\]

Donde $P_s$ representa la potencia de la señal y $P_n$
la potencia del ruido.

\begin{itemize}
    \item \textbf{30 dB:} el ruido es bajo y la detección
    de bits es más confiable.
    \item \textbf{10 dB:} aumenta la influencia del ruido
    y pueden aparecer errores de detección.
    \item \textbf{0 dB:} ambas potencias son iguales,
    provocando una mayor probabilidad de errores.
\end{itemize}

En consecuencia, una menor SNR deteriora la señal recibida
y afecta la calidad del audio reconstruido.

\subsection{Filtro pasabajas}

\textbf{¿Cómo cambia la señal al modificar la frecuencia
de corte a 2000, 3400 y 5000 Hz?}

El filtro pasabajas conserva las componentes inferiores
a su frecuencia de corte y atenúa las superiores.

Con 2000 Hz se pierden componentes importantes de la voz,
produciendo un sonido apagado. Con 3400 Hz se conserva
aproximadamente el ancho de banda telefónico. Con 5000 Hz
se mantienen más frecuencias altas, aunque también pueden
permanecer componentes no deseadas.

\subsection{Orden del filtro}

\textbf{¿Qué diferencias existen al utilizar órdenes
2, 4 y 10?}

El orden determina la pendiente de atenuación del filtro.
Un orden 2 presenta una transición más suave, mientras
que los órdenes 4 y 10 ofrecen una mayor selectividad
en frecuencia.

En esta práctica se utiliza \texttt{filtfilt()}, que
realiza el filtrado en ambas direcciones y evita el
retardo de fase habitual de un filtro causal. Sin embargo,
pueden aparecer efectos transitorios.

\subsection{Comparación temporal}

\textbf{¿Qué tan similares son las formas de onda al
comparar intervalos de 100 ms?}

Al aumentar el intervalo de visualización se observa una
mayor porción de las señales original y reconstruida.

Ambas conservan características generales similares,
aunque pueden existir diferencias de amplitud y distorsiones
producidas por la cuantificación, el ruido del canal y
el filtrado.

\subsection{Comparación perceptual}

\textbf{¿Cuál es el efecto del filtrado pasabajas sobre
la calidad percibida de la voz?}

Antes del filtrado pueden percibirse irregularidades y
componentes de alta frecuencia no deseadas. Después del
filtrado, el sonido suele ser más suave debido a la
atenuación de dichas componentes.

Sin embargo, una frecuencia de corte demasiado baja
puede eliminar información importante de la voz y
disminuir su claridad.

\subsection{Medida cuantitativa: MSE}

\textbf{¿Cómo se calcula el error cuadrático medio y cómo
se relaciona con la percepción auditiva?}

El error cuadrático medio permite medir la diferencia
entre la señal original y la reconstruida:

\[
MSE=\frac{1}{N}\sum_{i=1}^{N}
(x_i-\hat{x}_i)^2
\]

Donde $N$ representa el número de muestras, $x_i$ la
señal original y $\hat{x}_i$ la señal reconstruida.

En Python se calcula mediante:

\begin{lstlisting}[language=Python]
mse = np.mean((x_in - x_filt)**2)
print("MSE:", mse)
\end{lstlisting}

Un valor cercano a cero indica mayor similitud entre
ambas señales, mientras que un valor elevado representa
una mayor diferencia. Sin embargo, el MSE no garantiza
por sí solo una mejor calidad auditiva.

% CONCLUSIONES
\section{Conclusiones}

\begin{enumerate}
    \item La frecuencia de muestreo y la resolución de
    cuantificación son fundamentales para conservar
    la información de una señal analógica.

    \item Las técnicas A-law y $\mu$-law permiten mejorar
    la representación de señales de voz mediante
    cuantificación no uniforme.

    \item El ruido presente en el canal puede provocar
    errores durante la detección de bits y afectar
    directamente la calidad del audio reconstruido.

    \item El filtrado pasabajas permite atenuar componentes
    no deseadas y obtener una señal más suave cuando
    sus parámetros se seleccionan adecuadamente.

    \item El error cuadrático medio proporciona una medida
    cuantitativa que complementa la comparación visual
    y auditiva de las señales.
\end{enumerate}

En general, la práctica permitió comprender el funcionamiento
de un sistema de digitalización PCM y la importancia de
cada una de sus etapas para conseguir una transmisión
de voz confiable.

\end{document}